\documentclass{article}
\usepackage{graphicx} % Required for inserting images
\usepackage{amsmath}

\title{TFG Física: Dispersión de la luz por partículas esféricas y conglomerados de partículas esféricas}
\author{Gorki}
\date{2024}

\begin{document}

\maketitle

% \section{Introducción}
% En este trabajo se va a realizar un estudio sobre 
% la dispersión de la luz por partículas esféricas y 
% conglomerados de partículas esféricas, utilizando 
% simulaciones basadas en la teoría de Mie. Gustav Mie 
% (1868-1957), físico Alemán, formuló la teoría que 
% proporciona una descripción matemática precisa de 
% cómo la luz se dispersa y absorbe cuando incide sobre 
% partículas esféricas de diferentes tamaños y composiciones.
% En época posterior, Michenko escribió un programa en Fortran
% basado en la teoría de Mie, que permite calcular
% las propiedades de dispersión y absorción de partículas
% esféricas individuales y de conglomerados de partículas.


% El principal objetivo de este trabajo es estudiar el comportamiento
% de la luz al interactuar con partículas esféricas. Para lograr el objetivo
% principal, se plantean los siguientes objetivos específicos:

% \begin{itemize}
%     \item Analizar el programa de Fortran de Mishenko.
%     \item Caracterizar la dispersión de la luz por una única partícula esférica
%     \item Caracterizar la dispersión de la luz por múltiples partículas esféricas
%     \item Caso de estudio: Aplicación de la teoría de Mie en nanoantenas
% \end{itemize}

\section{Introducción}
\hspace{1em} En este trabajo se realizará un estudio sobre la dispersión de la 
luz por partículas esféricas y conglomerados de partículas 
esféricas, utilizando simulaciones basadas en la teoría de Mie. 
Formulada por el físico alemán Gustav Mie a principios del 
siglo XX, esta teoría ofrece una descripción matemática 
detallada de cómo la luz interactúa con partículas de 
diferentes tamaños y composiciones. A lo largo de los años, 
la teoría de Mie se ha consolidado como una herramienta 
fundamental para entender fenómenos ópticos en diversas 
disciplinas.

La importancia de estudiar la dispersión de la luz 
mediante la teoría de Mie radica en su aplicabilidad 
en una amplia gama de campos, como la óptica atmosférica, 
biomédica, y el desarrollo de materiales fotónicos. En 
estos contextos, la interacción de la luz con partículas 
esféricas de tamaño comparable a la longitud de onda de 
la luz puede influir de manera significativa en el 
comportamiento de la radiación, afectando tanto su 
dispersión como su absorción.

El objetivo principal de este trabajo es estudiar 
cómo se dispersa la luz cuando interactúa con partículas 
esféricas, utilizando simulaciones numéricas basadas 
en el programa de Fortran de D. W. Mackowski, que implementa 
la teoría de Mie. Para ello, se plantean los siguientes 
objetivos específicos:

\begin{itemize}
\item Analizar el programa de Fortran de D. W. Mackowski.
\item Caracterizar la dispersión de la luz por una única 
        partícula esférica.
\item Estudiar la dispersión de la luz por conglomerados 
        de partículas esféricas.
\item Caso de estudio: Aplicación de la teoría de Mie en 
        nanoantenas ópticas.
\end{itemize}

\section{Fundamento teórico}
\subsection{Dispersión y absorción}
%¿Qué es la dispersión de la luz? 
La dispersión de la luz ocurre cuando la radiación electromagnética, 
al chocar con partículas, cambia de dirección en su trayectoria de 
propagación. Este fenómeno es consecuencia de la excitación de partículas 
y la posterior radiación, y durante este proceso, parte de la energía 
puede perderse en forma de absorción.


Imaginemos una partícula dividida en pequeñas regiones. Cuando una 
onda electromagnética incide en la partícula, el campo eléctrico de 
la onda induce la polarización de la partícula, haciendo que las pequeñas 
regiones adquieran momentos dipolares cuya polaridad oscila con la misma 
frecuencia que la onda. Estos dipolos oscilantes actúan como fuentes 
emisoras de radiación electromagnética secundaria. El campo electromagnético 
detectado en un punto será la superposición de todas estas radiaciones secundarias. 
La magnitud y dirección de esta radiación total dependerán de la fase y 
amplitud de cada una de las radiaciones individuales, que a su vez están 
determinadas por las propiedades ópticas del material. Por lo tanto, la 
dispersión de la luz depende de las propiedades ópticas de la partícula.


En 1865, Maxwell publicó sus famosas ecuaciones, que serán 
necesarias para desarrollar la teoría de Mie. El conjunto 
de ecuaciones del electromagnetismo está asociado a cuatro 
vectores de campo: El campo eléctrico $\vec E$, el 
desplazamiento dieléctrico $\vec D$, la intensidad del 
campo magnético $\vec H$ y la inducción magnética $\vec B$. 
En cada punto del espacio, los vectores están sujetos a 
las ecuaciones de Maxwell y la versión macroscópica de 
estas ecuaciones es la siguiente:

\begin{equation}
    \vec{\nabla} \cdot \vec{D} = \rho_F
\end{equation}
\begin{equation}
    \vec{\nabla} \times \vec{E} + \frac{\partial \vec{B}}{\partial t} = 0
\end{equation}
\begin{equation}
    \vec{\nabla} \cdot \vec B = 0
\end{equation}
\begin{equation}
    \vec \nabla \times \vec H = \vec J_F + \frac{\partial\vec D}{\partial t}
\end{equation}

Aquí, $t$ es el tiempo, $\rho_F$ la densidad de carga libre y $\vec J_F$ la densidad de corriente. Las relaciones entre el campo eléctrico e inducción magnética están relacionadas con el desplazamiento y campo magnético por:

\begin{equation}
    \vec D = \epsilon_0 \vec E + \vec P
\end{equation}
\begin{equation}
    \vec H = \frac{\vec B}{\mu_0} - \vec M
\end{equation}

$\vec P$ y $\vec M$ son la polarización eléctrica y la magnetización, y $\epsilon_0$ y $\mu_0$ la permitividad y permeabilidad magnética en el vacío. Por último, las relaciones constitutivas:

\begin{equation}
    \vec J_F = \sigma \vec E
\end{equation}
\begin{equation}
    \vec B = \mu \vec H
\end{equation}
\begin{equation}
    \vec P = \epsilon_0\chi \vec E
\end{equation}

donde $\sigma$ es la conductividad, $\mu$ la permeabilidad y $\chi$ la susceptibilidad eléctrica. Los coeficientes son dependientes de la frecuencia y del medio. %, pero vamos a tomarlos como independientes del medio por ahora.



Un campo electromagnético en un medio lineal, homogéneo e isotrópico tiene que satisfacer la ecuación de onda:
\begin{equation}
    \nabla^2\vec E + k^2 \vec E = 0, \ \  \nabla^2\vec H + k^2 \vec H = 0
\end{equation}

donde $k^2 = \omega^2\varepsilon\mu $ y la divergencia de anula.

\begin{equation}
    \nabla \cdot \vec E = 0, \ \ \nabla \cdot \vec H = 0
\end{equation}

\begin{equation}
 \nabla \cdot \vec{E}=0, \quad \nabla \cdot \vec{H}=0 \\
 \nabla \times \vec{E}=i \omega \mu \vec{H}, \quad \nabla \times \vec{H}=-i \omega \epsilon \vec{E}
\end{equation}

Las ondas electromagnéticas planas tienen la siguiente forma:
\begin{equation}
    \vec E_c = \vec E_0 e^{i\vec k \cdot \vec x-i\omega t}; \ \ \vec H_c = \vec H_0 e^{i\vec k \cdot \vec x-i\omega t}
\end{equation}


La forma general de los campos es la siguiente:
\begin{equation}
    \vec F = \vec A \cos(\omega t) + \vec B \sin(\omega t)
\end{equation}

y la forma compleja es:
\begin{equation}
    \vec F_c = \vec C e^{-i\omega t} \ ; \ \vec C = \vec A + i\vec B.
\end{equation}

El exponente puede ser negativo o positivo, pero vamos a seguir la convención de ser negativo.

Es de alto interés definir el vector de Poynting, debido a que es de fundamental importancia en problemas de dispersión electromagnética. Se define como $\vec S = \vec E \times \vec H$, la magnitud representa la tasa de transferencia de energía y su dirección es la dirección de propagación.



La velocidad neta a la que la energía electromagnética atraviesa una área que encierra un volumen V es:
\begin{equation}
    W = -\int_A \vec S \cdot \hat n  \ dA
\end{equation}

Si W es positiva, entonces la energía electromagnética ha sido absorbida en el volumen, es decir, la energía ha sido convertida en otra forma de energía en V.

Como la mayoría de instrumentos no pueden calcular $\vec S$, se utiliza la media:

\begin{equation}
    \langle \vec{S} \rangle = \frac{1}{2} \, \text{Re}\!\left\{ \vec{E}_c \times \vec{H}^*_c \right\}.
\end{equation}

Las ondas electromagnéticas planas:
\begin{equation}
    \vec E_c = \vec E_0 e^{i\vec k \cdot \vec x-i\omega t}; \ \ \vec H_c = \vec H_0 e^{i\vec k \cdot \vec x-i\omega t}
\end{equation}

El vector de onda de una onda homogénea

\begin{equation}
    k = k' + ik'' = \frac{\omega N}{c},
\end{equation}
donde N es el índice de refracción complejo:
\begin{equation}
    N = c\sqrt{\varepsilon\mu} = \sqrt{\frac{\varepsilon\mu}{\varepsilon_0 \mu_0}} = n + ik
\end{equation}

Los números $n$ y $k$ se conocen como constantes ópticas (que no son constantes, ya que dependen de la frecuencia). La parte real del índice de refracción determina la velocidad a la que se propaga la onda en el medio, mientras que la parte imaginaria, determina la atenuación de la onda a medida que se propaga por el medio.

El número de onda $k = \omega/v$ donde $v$ es la velocidad de propagación de la onda en el medio.
La longitud de onda en el medio es $\lambda = 2\pi/k$.



La polarización es una propiedad intrínseca de las ondas electromagnéticas monocromáticas. La parte real del campo eléctrico viene dado por

\begin{equation}
    \vec E = \vec A \cos{kz - \omega t} - \vec B \sin{kz -\omega t}
\end{equation}

que, al seleccionar un plano arbitrario, como $z = 0$, el campo eléctrico sería

\begin{equation}
    \vec E(z=0) = \vec A \cos \omega t + \vec B \sin \omega t
\end{equation}

que describe una elipse. Si $\vec A$ o $\vec B$ es 0, la vibración de la elipse sería solo una línea recta, y la onda estaría polarizada linealmente.

Para hacer una descripción precisa de una onda monocromática, se utilizan los parámetros de Stokes. Los parámetros de Stokes [Stokes (1819-1903)] son un grupo de valores que describen el estado de polarización de una onda electromagnética.

\begin{equation*}
    I = c^2
\end{equation*}

\begin{equation*}
    Q = c^2 \cos 2\eta \cos 2\gamma
\end{equation*}
\begin{equation*}
    U = c^2 \cos 2\eta  \sin 2\gamma
\end{equation*}
\begin{equation}
    V = c^2 \sin2\eta
\end{equation}
donde $c^2 = a^2+b^2$ 

I es la irradiancia de la onda, que es la magnitud utilizada para describir la potencia incidente por unidad de superficie de radiación electromagnética.

AQUI FALTA PONER LA IMAGEN DE QUE ES ETA Y GAMMA, son los parametros de elipcicidad.

Los parámetros de Stokes están definidos en términos medibles, (irradiancia). Los parámetros no son independientes, tenemos las siguientes relaciones:

\begin{equation}
    \tan 2\gamma = \frac{U}{Q}
\end{equation}

\begin{equation}
    I^2 = Q^2+U^2+V^2
\end{equation}

La matriz de Mueller describe la relación entre el vector de Stokes de la onda incidente y el vector de onda de la onda transmitida.

por ahora no lo voy a explicar porque no parece esencial.




Extinción dispersión y absorción:
Si construimos una esfera imaginaria alrededor de una partícula, la velocidad neta a la que la energía electromagnética atraviesa la superficie A de la esfera es

\begin{equation}
    W_a = -\int_A \vec S \cdot \hat{e}_rdA
\end{equation}

Si $W_a$ > 0 la energía es absorbida dentro de la esfera y como el medio es no absorbente, la energía tiene que haber sido absorbida por la partícula.

Como $\vec S = \vec S_i + \vec S_ext + S_{ext}$, entonces
\begin{equation}
    W_a = W_i -W_s + W_{ext}
\end{equation}

$W_s$ es la velocidad de a la que la energía se dispersa a través de la superficie.

\begin{equation}
    W_i = -\int_A \vec S_i \cdot \hat{e}_rdA, \ \
    W_s = \int_A \vec S_s \cdot \hat{e}_rdA, \ \
    W_{ext}= -\int_A \vec S_{ext} \cdot \hat{e}_rdA.
\end{equation}

como $W_i$ se desvanece en un medio no absorbente:

\begin{equation}
    W_ext = W_a + W_s
\end{equation}

Las secciones eficaces son 
\begin{equation}
    C_{ext} = C_{abs} + C_{sca}
\end{equation}

Donde $C_{abs}=W_{abs}/I_i$ y $C_{sca}=W_{sca}/I_i$ 


p es $|\vec X|^2/k^2C_{sca}$ y está normalizado en el ángulo sólido.

La media del ángulo del coseno de dispersión, el parámetro de asimetría g es
\begin{equation}
    g = <\cos\theta> = \int_{4\pi} p\cos \theta
\end{equation}

Las eficiencias para la dispersión extinción y absorción son:

\begin{equation}
    Q_{sca} = \frac{C_{Sca}}{G} \ \ Q_{ext} = \frac{C_{ext}}{G} \ \ Q_{abs} = \frac{C_{abs}}{G}
\end{equation}

G es el área de sección transversal perpendicular a la onda incidente.

Hasta ahora hemos considerado una única partícula. Con múltiples partículas las secciones eficaces serian:

\begin{equation}
    C_{ext} = \sum_j C_{ext, j}; \ \ C_{ext, j} = C_{abs,j} + C_{sca, j}
\end{equation}


Absorción y dispersión por una esfera.

\section{Referencias}
\bibliographystyle{plain}
\bibliography{referencias}
\end{document}
